%mac

\setcounter{paragraph}{0}
\paragraph{Lực cản của chất lưu}
\immini{\begin{itemize}
\item Thông thường thuật ngữ chất lưu được dùng để chỉ chất lỏng và chất khí.
\item Mọi vật chuyển động trong chất lưu luôn chịu tác dụng bởi lực cản của chất lưu.
\item Lực này ngược hướng với chuyển động và cản trở chuyển động của vật.
\item Lực cản của chất lưu có tác dụng tương tự như lực ma sát, chúng làm chuyển động của các vật bị chậm lại.
\item Lực cản phụ thuộc vào hình dạng và tốc độ của vật.
\end{itemize}}{\begin{tikzpicture}
\node[] at (0,0){\includegraphics[width=5cm]{images/19-1}};
\end{tikzpicture}}	
\paragraph{Lực nâng của chất lưu}
\begin{itemize}
\item Khi vật chuyển động trong nước hay trong không khí thì ngoài lực cản (của không khí, của nước), vật còn chịu tác dụng của lực nâng.
\end{itemize}
\immini{	\begin{itemize}
\item Hình bên cho thấy, nếu cánh máy bay nghiêng đi một chút, chếch theo luồng gió, thì tăng thêm lực nâng tác dụng vuông góc với mặt dưới của cánh để nâng máy bay lên cao.
\item Lực nâng của chất lưu giúp khinh khí cầu lơ lửng trên không trung, máy bay di chuyển trong không khí, cho phép tàu thuyền di chuyển trên mặt nước,...
\end{itemize}}{\begin{tikzpicture}
\node[] at (0,0){\includegraphics[width=5cm]{images/19-3}};
\end{tikzpicture}}
\immini{\begin{itemize}
\item Nhờ có lực nâng
\begin{itemize}
\item Máy bay có thể di chuyển trong không khí.
\item Tàu thuyền có thể nổi và di chuyển được trên mặt nước.
\item Khinh khí cầu lơ lửng trên không trung.
\item  Nhiều sinh vật bay lượn dễ dàng trong không khí.
\end{itemize}	
\end{itemize}}{\begin{tikzpicture}
\node[] at (0,0){\includegraphics[width=5cm]{images/19-6}};
\end{tikzpicture}
}
\begin{center}
\begin{tikzpicture}
\node[] at (0,0){\includegraphics[width=5cm]{images/19-4}};
\node[] at (7,0){
\includegraphics[width=5cm]{images/19-5}};
\end{tikzpicture}
\end{center}	
\Note{Khi vật rơi trong chất lưu dưới tác dụng của trọng lực và lực cản của chất lưu thì đến một lúc nào đó vật sẽ đạt tới vận tốc giới hạn và sẽ chuyển động với vận tốc này}
\paragraph{Lực đẩy Archimedes}
\begin{itemize}
\item Lực đẩy Archimedes là trường hợp riêng của lực nâng vật đứng yên trong chất lưu.
\item Công thức tính lực đẩy Archimedes: 
\begin{align*}
\boder{F_A=\rho g V}
\end{align*}
Trong đó:
\begin{itemize}
\item $F_A$: lực đẩy Archimedes (N).
\item $\rho$: Khối lượng riêng của chất lỏng $(kg/m^3)$
\item V: thể tích phần chất lỏng bị chiếm chỗ $(m^3)$
\end{itemize}
\end{itemize}	



